\chapter{Parent-селектор кратности и скорости коннектора c0}
\label{ch:v8-baryon-c0-connector-multiplicity-and-rate-parent-selector-gate}

Допустимый старо-новый коннектор задаётся вещественным вектором
\begin{equation}
 z=(z_0,z_1,z_2)\in\mathbb R^3,
 \qquad \lVert V(z)\rVert_{\rm HS}^2=z_0^2+z_1^2+z_2^2.
 \label{eq:v8-baryon-c0-connector-multiplicity-trace-norm}
\end{equation}
Три координаты принадлежат различным концевым бимодульным меткам, но все
имеют один gauge-тип $(1,1)_{-1}$. Текущий parent определён на старом
42-кадре и не содержит старо-новых координат, поэтому его нулевое
продолжение даёт
\begin{equation}
 \Hess S_{\rm old}\big|_{\mathcal V_{-1}^{\rm odd}}=0_3.
 \label{eq:v8-baryon-c0-old-parent-zero-connector-hessian}
\end{equation}

\section{Изотропный общий след не выбирает точку RP2}

Наиболее экономный общий след допускает только радиальный потенциал
\begin{equation}
 S_{\rm iso}(z)=\frac a2\lVert z\rVert^2
 +\frac b4\lVert z\rVert^4,
 \qquad b>0.
 \label{eq:v8-baryon-c0-isotropic-connector-parent}
\end{equation}
При $a<0$ его ненулевые стационарные точки образуют сферу, а после
отождествления $z\sim-z$ --- всё прежнее семейство:
\begin{equation}
 \lVert z\rVert^2=-\frac ab,
 \qquad \mathcal M_{\rm vac}=S^2/\{\pm1\}=\mathbb{RP}^2.
 \label{eq:v8-baryon-c0-isotropic-vacuum-rp2}
\end{equation}
В точке $z_*=r e_0$, где $a=-br^2$, гессиан равен
\begin{equation}
 \Hess S_{\rm iso}(z_*)=2br^2 e_0e_0^{\mathsf T},
 \qquad \rank\Hess S_{\rm iso}(z_*)=1,
 \qquad \dim\ker=2.
 \label{eq:v8-baryon-c0-isotropic-parent-two-flat-directions}
\end{equation}
Следовательно, общий trace фиксирует норму, но не направление коннектора.

Даже гауссовская флуктуационная модель в симметричной фазе даёт лишь
\begin{equation}
 C_{\rm conn}=T(m^2I_3)^{-1}=\frac{T}{m^2}I_3,
 \qquad \gamma\propto\frac{T}{m^2},
 \label{eq:v8-baryon-c0-isotropic-gaussian-rate}
\end{equation}
где отношение $T/m^2$ не выводится из следовой метрики.

\section{Что потребовалось бы для выбора ветви}

Полные концевые проекторы разрешают анизотропный член
\begin{equation}
 S_{\rm lab}^{(2)}(z)=\frac12
 \left(m_0^2z_0^2+m_1^2z_1^2+m_2^2z_2^2\right),
 \label{eq:v8-baryon-c0-labeled-connector-quadratic-parent}
\end{equation}
но три коэффициента являются новыми parent-данными. Например, два
одинаково допустимых выбора
\begin{equation}
 M_A=\operatorname{diag}(-1,1,2),
 \qquad M_B=\operatorname{diag}(2,-1,1)
 \label{eq:v8-baryon-c0-two-admissible-branch-selectors}
\end{equation}
выбирают разные координатные ветви и сохраняют gauge-, grading- и
Real-условия. Симметрия не различает эти варианты.

Наконец, сам генератор имеет точную перенормировочную свободу
\begin{equation}
 V\mapsto sV,
 \qquad \gamma\mapsto\frac{\gamma}{s^2},
 \qquad \gamma\,\mathcal D[V]
 =\frac{\gamma}{s^2}\mathcal D[sV].
 \label{eq:v8-baryon-c0-connector-rate-rescaling-orbit}
\end{equation}
Даже после фиксации $\lVert V\rVert_{\rm HS}=1$ абсолютная $\gamma$
остаётся интенсивностью среды, отсутствующей в конечном parent action.

Итоговый реестр селектора равен
\begin{equation}
 N_{\rm derived}=0/2:
 \qquad [z]\in\mathbb{RP}^2\ \text{не выбран},
 \qquad \gamma>0\ \text{не выбрана}.
 \label{eq:v8-baryon-c0-parent-selector-zero-of-two}
\end{equation}

\begin{theorem}[Радиальная полнота и угловой дефицит]
Общий след канонически задаёт евклидову норму на трёхмерном пространстве
коннекторов. Любой parent, зависящий только от этой нормы, оставляет
двумерную угловую моду и не выбирает точку $\mathbb{RP}^2$.
\end{theorem}

\begin{nogocriterion}[Кратность и скорость требуют разных новых данных]
Бимодульно-анизотропный гессиан может выбрать направление, но его
коэффициенты не следуют из текущего parent. Абсолютная скорость дополнительно
требует состояния или интенсивности среды. Нельзя использовать один
произвольный коэффициент одновременно как селектор ветви и как физические
часы процесса.
\end{nogocriterion}

\begin{protocol}\begin{enumerate}
\item размерность multiplicity-пространства: \textbf{$3$};
\item метрика общего следа: \textbf{$I_3$};
\item текущий parent Hessian на новых координатах: \textbf{$0_3$};
\item изотропный ненулевой vacuum: \textbf{$\mathbb{RP}^2$};
\item угловых нулевых мод: \textbf{$2$};
\item анизотропный селектор допустим: \textbf{да, с новыми весами};
\item направление коннектора выведено: \textbf{нет};
\item абсолютная скорость выведена: \textbf{нет};
\item реестр: \textbf{$0/2$};
\item следующий гейт: \textbf{происхождение весов из расширенной
концевой бимодульной алгебры}.
\end{enumerate}\end{protocol}

Машинный сертификат сохранён в
\path{s2t_v8_baryon_c0_connector_multiplicity_and_rate_parent_selector_gate_results.json}.